\section*{Overview}

Lowest common ancestor is the standard query on rooted trees: for two vertices \(u\) and \(v\), find the deepest
vertex that is an ancestor of both. Binary lifting is the contest workhorse because it is easy to combine with depth
logic and other ancestor-style queries.

\section*{When to Use It}

Use LCA when:

\begin{itemize}
  \item a problem asks about paths in a rooted tree,
  \item distance, ancestry, or path decomposition repeatedly mentions two vertices,
  \item you need a reusable tree-query baseline before heavier tools like HLD.
\end{itemize}

\section*{Core Idea}

Precompute \(\texttt{up[k][v]}\), the \(2^k\)-th ancestor of vertex \(v\). Then:

\begin{itemize}
  \item raise the deeper node until both nodes are at the same depth,
  \item lift both nodes together from large powers to small until their ancestors diverge,
  \item the parent just above that divergence is the LCA.
\end{itemize}

\section*{Key Insight}

Binary lifting turns ancestor climbing into bit decomposition. Moving up \(13\) levels is the same as moving up
\(8 + 4 + 1\) levels. That is exactly what the jump table stores.

\section*{Operations / Main Technique}

\begin{itemize}
  \item DFS or BFS to compute parent and depth,
  \item build the jump table,
  \item answer \texttt{lca(u, v)},
  \item support \texttt{kth\_ancestor(v, k)} and distance queries as well.
\end{itemize}

\paragraph{Worked example.}
If one node is 11 levels deeper than the other, binary lifting raises it by the set bits of \(11\), namely
\(8 + 2 + 1\), before the paired lifting phase starts.

\section*{Correctness Intuition}

The first phase preserves the LCA because only the deeper node moves upward. In the second phase, if the
\(2^k\)-ancestors of the two nodes differ, both nodes can safely jump there without passing the LCA. After all such
jumps, they sit directly below the LCA.

\section*{Complexity Analysis}

\begin{itemize}
  \item preprocessing: \(O(n \log n)\),
  \item one LCA query: \(O(\log n)\),
  \item memory: \(O(n \log n)\).
\end{itemize}

\section*{Implementation}

The code builds the binary-lifting table from a rooted tree and supports:

\begin{itemize}
  \item \texttt{kth\_ancestor},
  \item \texttt{lca},
  \item \texttt{distance}.
\end{itemize}

\section*{Common Pitfalls}

\begin{itemize}
  \item Forgetting to root the tree consistently.
  \item Using too few lifting levels for the maximum \(n\).
  \item Mixing edge count and vertex depth conventions in distance formulas.
  \item Recursion depth issues on deep trees if DFS is used directly.
\end{itemize}

\section*{Variants / Extensions}

\begin{itemize}
  \item Euler tour plus RMQ for \(O(1)\) LCA queries after different preprocessing.
  \item Binary lifting for functional graphs.
  \item Path queries combined with HLD or Euler-tour flattening.
\end{itemize}

\section*{Practice Problems}

\begin{itemize}
  \item Distance queries on a tree.
  \item K-th ancestor or jump queries.
  \item Path intersection or meeting-point problems on trees.
\end{itemize}

\section*{References}

\begin{itemize}
  \item \href{https://cp-algorithms.com/graph/lca_binary_lifting.html}{cp-algorithms: Lowest Common Ancestor - Binary Lifting}.
  \item \href{https://usaco.guide/plat/binary-jump}{USACO Guide: Binary Jumping}.
  \item \href{https://codeforces.com/catalog}{Codeforces Catalog}.
\end{itemize}
